\appendix
\section{Exponential Tailbounds}

I just took these from my notes in random matrices.

\begin{lemma}
  \label{lemma:hoeff}
  (Hoeffding's Lemma) Let $X$ be a scalar r.v taking values in
  $[a,b]$. Then for any $t > 0$
  \begin{equation}
    \label{lemma:hoef:eq1}
    \E e^{tX} \leq e^{t\Ex X}(1 + O(t^2\sigma^2\exp(O(t(b-a))))).
  \end{equation}
  In particular
  \begin{equation}
    \label{lemma:hoef:eq2}
    \E e^{tX} \leq e^{t\E X}\exp(O(t^2(b-a)^2)).
  \end{equation}
\end{lemma}

\begin{proof}
  Normalize $X$ as to $\Ex(X) = 0$ and $b - a = 1$. And consider the
  Taylor expansion of $e^{tX}$,
  \[
    e^{tX} = 1 + tX + O(t^2X^2\exp(O(t)))
  \]
  taking expectation yields the first result. The second one follows
  from assymptotics of the exponential function.
\end{proof}

\begin{theorem}
  \label{trm:McDiamird}
  (McDiamird's Theorem) Let $(R_i)_{i=1}^n$ be a tuple of ranges
  and $X = (X_i)_{i=1}^{n}$ a random vector taking values in
  their product.
  Let $F: \prod_{i=1}^n R_i \to \R$ be a coordinate-wise bounded
  function (this is a strong Lipchitz condition). That is changing
  any coordinate $x_i$ may change at most $c_i$ the total value.
  Then
  \[
    \Pr( |F(X) - \Ex F(X)| > \lambda \sigma) \leq C \exp(-c \lambda^2)
  \]
  for some $C,c > 0$ and $\sigma^2 = \sum_{i=1}^n c_i^2$.
\end{theorem}

\begin{proof}
  By simmetry, we may only consider $\Pr(F(X) - \E F(X) > \lambda
  \sigma)$. We turn to Hoeffding's Lemma [\ref{lemma:hoeff}],
  consider
  \[
    \Ex \exp(tF(X)) = \Ex \big[ \Ex[\exp(tF(X))| X_1, \dots X_{n-1}] \big],
  \]
  we should try to understand $\Ex[\exp(tF(X))| X_1, \dots X_{n-1}]$.
  We have that
  \[
    \Ex[\exp(tF(X))| X_1, \dots X_{n-1}] = \Ex[\exp(t(F(X) -
    \Ex(F)))\exp(t\Ex F(X))| X_1, \dots X_{n-1}]
  \]
  applying (\ref{lemma:hoef:eq2}) to the first product, (here we are
  conditioning on $X_1, \dots X_n$)
  \[
    \Ex[\exp(t(F(X) - \Ex(F)))\exp(t\Ex F(X))| - ] \leq
    \exp(O(t^2c_n^2))\Ex[\exp(t\Ex F(X)) | -]
  \]
  and we can iteratively apply the theorem to the new coordinate-wise
  bounded function
  \[
    G(X_1, \dots X_{n-1}) = \Ex[F(X)|X_1, \dots X_{n-1}].
  \]
  Notice $G$ has the same coefficients as $F$, because
  \begin{align*}
    |G(y, x_2 \dots, x_{n-1}) - G(z, x_2 \dots, x_{n-1})| &= \bigg|
    \int_{R_n} F(y,\dots,x_n) dx_n - \int_{R_n} F(z,\dots,x_n) dx_n\bigg|\\
    &= \bigg| \int_{R_n} F(y,\dots,x_n) - F(z,\dots,x_n) dx_n\bigg| < c_1,
  \end{align*}
  and here we can take any other coordinate instead of the first.
\end{proof}

\begin{corollary}
  \label{trm:hoeffding_ineq}
  (Hoeffding's inequality) Let $(X_i)_{i=1}^n$ be real r.v's taking
  value in bounded intervals $[a_i,b_i]$. Then, if $S = \sum_{i \leq
  n} X_i$, we have
  \[
    \Pr( |S  - \Ex S| > \lambda\sigma) < C\exp(-c \lambda^2)
  \]
  for some $C,c > 0$, where $\sigma^2 = \sum_{i \leq n} (b_i-a_i)^2$.
\end{corollary}
\begin{proof}
  Apply [\ref{trm:McDiamird}] with $F(X_1, \dots X_n) = X_1 + \dots + X_n$.
\end{proof}

\begin{corollary}
  \label{lemma:expec_binomial_bounds}
  (Tailbounds for the Binomial distribution) Let $B(n,p)$ be the
  binomial distribution.
  We have that
  \[
    \Pr(|B(n,p) - np| \geq \delta np) \leq C\exp(-c \delta^2 (np)^2)
  \]
  for some absolute $C,c$.
\end{corollary}
\begin{proof}
  Note that
  \[
    \Pr(|B(n,p) - np| \geq \lambda \sqrt{np}) \leq \Pr(|B(n,p) - np|
    \geq \lambda \sqrt{np(1-p)}) \leq C\exp(-c\lambda^2)
  \]
  by \ref{trm:hoeffding_ineq}. Applying it with $\lambda = \delta
  \sqrt{np}$ yields the result.
\end{proof}

\begin{theorem}
  (Azuma's Inequality) Let $X_1, X_2, \dots X_n$ be a sequence scalar
  r.v's with $|X_i| \leq 1$ almost surely. Assume also
  that we have the martingale difference property
  \[
    \E(X_i | X_1, \dots, X_{i-1}) = 0
  \]
  almost surely for all $i = 1,\dots,n$. Then for any $\lambda > 0$,
  the sum $S_n := X_1 + \dots + X_n$ satisfies
  \[
    \Pr(|S_n| \geq \lambda \sqrt{n}) \leq C\exp(-c\lambda^2)
  \]
  for some absolute $C,c > 0$.
\end{theorem}

\begin{proof}
  We prove this by calculating exponential moments and using
  Hoeffding's Lemma again. Note that $S_n = S_{n-1} + X_n$ so that
  \[
    \Ex(e^{tS_n}) = \Ex(e^{tS_{n-1}}e^{tX_n}).
  \]
  By taking the expectation of conditioning on the first $n-1$
  values, we get that
  \[
    \Ex(e^{tS_n}) = \Ex \Ex (e^{tS_{n-1}}e^{X_n} | X_1, \dots,
    X_{n-1}) =\Ex [e^{tS_{n-1}}\Ex(e^{tX_n} | X_1, \dots, X_{n-1})]
  \]
  Now, by \ref{lemma:hoeff}, we have that $\Ex(e^{tX_n} | X_1, \dots,
  X_{n-1}) \leq e^{O(t^2)}$, so by iterating we get the result.

\end{proof}
